\section{Testing responses to a five-course menu}\label{sec:behavior-tests}
\paragraph{What the field study identifies}
\citet{smartplatforms} study a warning shown when a family tried to submit a list with high predicted nonplacement risk in 2018--20. The message encouraged adding schools without naming one. Tables II and III report the discontinuity at 30\% risk. The probability of adding any school just below the threshold was .199; crossing it raised that probability by .216 (SE .010). The rise in appending was .205 (.009). The estimates refer to overlapping actions because a family could add multiple schools. The probability of placement at a listed school rose by .038 (.009). In their sample, the treatment induced 1.573 extra schools per induced adder (.090). We verified the coefficients against the archived tables \citep{warningdata}. The assignment outcome for a warned family and the national placement difference between two recommenders are different estimands. In a separate Chilean experiment, access to an individualized report with suggested schools also changed list revision behavior \citep{schoolSearch}. Neither study observes our two menus.

\paragraph{National simulations}
We use the two frozen menus, the original quota, the same four complete national cohorts, and the reconstructed allocation from Section~\ref{sec:chile}. An addition can come from a course the family ranked on its original eligible report. In the school-level tests, its establishment must be absent from all courses on the displayed menu. The first candidate is the highest-ranked qualifying course. The alternative is the one with the lowest ratio of eligible reports plus one to published vacancies plus one. Submitted rank breaks ties. This proxy uses no enrollment or priority draw. It does not guarantee a seat. A course selected for a pupil is identical across the two ways of inserting it into that pupil's menu. We either insert it in its original rank or append it after the shown courses. Family promotions, sibling priorities, separate seat types and guaranteed returns are then processed afresh. A protected return can make a tail option unreachable; it still counts as a pursued school in the coverage calculation. A family with no qualifying option cannot add one in this test.

A common uniform draw for each pupil determines participation under both policies and across all rates. Intermediate rates use eight independently seeded repetitions in each cohort. Tables report the mean of the four-year sums. The variation across seeds measures uncertainty from this simulation design. At a response rate of .415, appending the first school gives 1,656 fewer linked placements on average (Monte Carlo SE 1.56); inserting it in submitted order gives 1,340.75 fewer (SE 4.01). To vary the number of additions, 57.3\% of the respondents in the .415 scenario try to add a second distinct omitted school. This is a modeling choice motivated by the field study's estimate of 1.573 schools per induced adder. The actual average is 1.289 under either menu because many reports contain only one recoverable school. Tail insertion gives 2,817.6 more April hits and 1,645 fewer linked placements. Original-order insertion gives the same hits and 1,279 fewer linked placements. We also tested probabilities .10, .216 and .50. Under the append rule we crossed all seven tested response rates (including zero and one) for the two menus. Linked placement fell in every draw at all 49 rate pairs, with differences from $-1{,}694$ to $-1{,}608$. Coverage can change sign when uptake differs by policy. If everyone adds the first available school under the blind menu and no one does under the informed menu, the informed rule has 1,386 fewer hits and 1,694 fewer linked placements.

For a targeted check, we estimated nonplacement risk on the two existing displays with 64 simulated course-specific family lotteries per cohort, before using the recorded order for the final allocations. The stochastic check uses 2022--24. In those years 50,571 applicants had estimated risk between 20\% and 40\% on both displays. Only 5,907 could add a previously unshown school on their blind menu and 5,917 on their informed menu. With .415 participation restricted to that group, the mean linked placement difference was $-1{,}359.6$ if they appended it, compared with $-1{,}360$ under binding menus. In a separate check, giving the observed same-grade returners the least-congested candidate to append gives 4,671 more April hits and 1,625 fewer linked placements. The simulated threshold group is selected by a model of the frozen displays; the threshold in the field study used a different forecast.

\paragraph{Which families respond?}
Let $G_A$ and $G_B$ be the linked pupils whose blind or informed menu, respectively, misses their historical April school and would gain that school by adding its chosen candidate. Let $S$ be any common set of pupils willing to expand. Write $\Delta_0$ for informed minus blind coverage before anyone expands. Then
\[
  \Delta(S)=\Delta_0+|S\cap(G_B\setminus G_A)|-|S\cap(G_A\setminus G_B)|.
\]
This follows by summing the change in each pupil's hit indicator. Membership in both gain sets cancels. The minimum is $\Delta_0-|G_A\setminus G_B|$ and the maximum is $\Delta_0+|G_B\setminus G_A|$. Either endpoint is achieved by taking the corresponding exclusive set. More generally, among $U$ eligible pupils let $b$ and $r$ be the blind-only and informed-only gain counts. For exactly $m$ willing pupils the sharp lower endpoint is $\Delta_0-\min(m,b)+\max(0,m-U+r)$. The upper endpoint is $\Delta_0+\min(m,r)-\max(0,m-U+b)$. Sorting the individual contributions $-1,0,1$ and taking the first or last $m$ proves both claims. For the first omitted school the two exclusive counts are 4,817 blind and 774 informed, with 1,207 shared. For the least-congested option they are 3,718 blind and 1,256 informed, with 995 shared. The endpoints again hold when willingness is fixed at 19.9\% of applicants able to add under either menu, rounded down in each cohort. We supplement each exclusive group with pupils who are in neither exclusive group. This gives 34,742 willing applicant-years in each construction, out of 174,594 eligible. Enrollment labels enter these deliberate worst- and best-case constructions. They never enter the predictors or the ordinary response simulations.

With policy-specific marginal response probabilities $q_A,q_B$ common among each policy's potential gainers, expected coverage for the first-school rule is $4{,}638+1{,}981q_B-6{,}024q_A$. No independence across applicants is needed for this expectation. Under equal probabilities the expected gap is positive even at full response. For placements we clear the entire market for each concrete responder set. These sharp coverage bounds do not bound the simulated placement gap.

\paragraph{Complementary-round choices and seats}
The repeat-course comparison links each same-grade B2 applicant to their B1 record, excludes automatic continuity entries, and matches every explicitly ranked C2 school/course to the eligible original report. It adds every match that the pupil's policy omitted. The original-rank version yields 4,733 additional hits and 1,651 fewer linked placements. Appending in the complementary order yields the same hits and 1,653 fewer placements. Restricting the additions to the group with no historical regular offer through the wait list leaves the original linked placement gap of 1,670 under either insertion order. Among the 2,635 pupils who lose a binding-menu place in the full cohort, 408 subsequently reapplied at the same grade in the historical data, 390 explicitly named a new school and 122 subsequently enrolled at a newly listed school that had offered them an explicitly requested course. These counts describe their observed historical paths.

For the separate second-stage calculation let $U_A,U_B$ be the pupils unplaced under blind and informed first-stage menus. For a set $U$ define $r(U)$ as the maximum number that can be matched to the recorded second-stage schools using common capacities. Construct a unit-capacity edge from a source to each eligible pupil and from that pupil to each explicitly applied-to course; each course has an edge to the sink of capacity equal to its vacancies. Integral maximum flow gives the count. A minimum cut certifies each value. No first-stage recipient moves. For the linked estimates the graph uses only linked pupils and therefore gives them first claim on second-stage seats. Both policies use the same priority-free objective and complementary inventory. Later-only applicants and the effects of earlier hypothetical placements on released vacancies are omitted.

Even with other competitors in common, a useful bound holds for maximizing headcount. Restrict any matching on $U_B$ to $U_A\cap U_B$. It leaves a feasible matching on $U_A$, and the removed edges can cover at most $r(U_B\setminus U_A)$ pupils. Hence $r(U_B)-r(U_A)\leq r(U_B\setminus U_A)$. The same proof allows any set of extra common applicants with fixed options. On the new-school graph with published capacities this upper bound is 388 for all applicants and 358 for the linked-only graph. Closing the original deficits of 1,834 or 1,670 therefore requires some additional change to the graph, supply or objective. This inequality compares maxima, conditional on who can pursue which seats; two arbitrary priority-based assignments need not satisfy it.

\begin{table}[ht]\centering\small
\caption{Rank-informed minus rank-blind after a hypothetical second chance. Each arm keeps its initial recipients and maximizes placements among its unplaced pupils on observed second-round choices. Linked placement gives linked pupils first claim. A pursued hit includes a second-round choice only for a pupil unplaced by that policy and only at a course with positive common capacity.}\label{tab:second-chances}
\begin{tabular}{@{}llrrr@{}}\toprule
Courses allowed & Common supply & All placed & Linked placed & Pursued hits\\\midrule
New school & Published second round & $-1,558$ & $-1,419$ & $+4,752$\\
All explicitly listed & Published second round & $-1,550$ & $-1,409$ & $+4,763$\\
New school & Also cap at both residuals & $-1,580$ & $-1,439$ & $+4,743$\\
All explicitly listed & Also cap at both residuals & $-1,574$ & $-1,433$ & $+4,748$\\\bottomrule
\end{tabular}
\end{table}
The residual supply check caps each course by its available regular-round seats after each of the two initial allocations as well as its published complementary vacancies. A course absent from the regular inventory receives zero in this check. All four rows relax category restrictions and competition from other second-round applicants. The relation to an actual two-round allocation depends on which families would reapply, list new schools, release old places and accept a subsequent offer.

\paragraph{Why adding a backup can leave fewer pupils placed}
With fixed priorities, appending options can only increase a program's DA occupancy. Resume the old run with the extra proposals. A program either fills a vacancy or replaces its current lowest-priority holder. Chile's higher-grade sibling priority changes this argument. For any $k\ge 2$, appending one course for one applicant can reduce total placement by $k-1$.

To see this, suppose $H$ and $S$ each list only $u$ in the senior grade. There is one seat at $u$, which admits $H$, and a vacant senior seat at $x$. In the junior grade $k$ pupils list only $x$, and each of $S$'s $k$ younger siblings ranks $x$ above $y$. Both junior courses have $k$ seats. In the ordinary priority order at $x$, the pupils listing only $x$ outrank every sibling. They get $x$ and the siblings get $y$. Now append senior $x$ to $S$'s list. $S$ obtains that previously empty seat, giving all the younger siblings priority at junior $x$. They take its $k$ seats. The pupils listing only $x$ become unplaced and nobody takes junior $y$. The senior grade gains one recipient and the junior grade loses $k$. This example uses only regular seats, with no family preference promotion. The institutional solver reproduces it already at $k=2$, with six applicants and a fall from five placed to four.


\paragraph{New York and the pooled Chilean benchmark}
\begin{table}[ht]\centering\small
\caption{How the assumed response changes the comparison between recommenders. Entries are rank-informed minus rank-blind in Chile, and Haiku minus uniform in New York. Chilean totals use the fixed-priority pooled model over four application years. NYC uses all renters with the original geographic weights. In the last row of each panel both policies list the same eligible options.}\label{tab:response}
\begin{tabular}{@{}lrrr@{}}\toprule
 & Placements & First choices & Past offers on list\\\midrule
\multicolumn{4}{l}{\textit{Chile, changes in pupils}}\\
Pursue the five shown & $-1,787$ & $+2,218$ & $+5,552$\\
Restore one at its original rank & $-940$ & $-846$ & $+676$\\
Append that same option last & $-1,761$ & $+2,163$ & $+676$\\
Append every omitted option & $-1,331$ & $+2,005$ & $0$\\\midrule
 & Ordinary placed & Enriched placed & Enriched first choices\\\midrule
\multicolumn{4}{l}{\textit{NYC, changes in percentage points}}\\
Pursue the five shown & $-2.71$ & $-5.40$ & $+2.87$\\
Restore one at its original rank & $-1.98$ & $-4.28$ & $+0.18$\\
Append that same option last & $-2.17$ & $-4.41$ & $+2.97$\\
Append every omitted option & $-0.38$ & $+0.28$ & $+1.84$\\\bottomrule
\end{tabular}
\end{table}

In this earlier pooled comparison, to restore one option, take each applicant's highest originally ranked eligible option missing from that policy's menu, if one exists. For a given policy the recovered option is the same in both insertion scenarios. This change alone flips the first-choice comparison in \emph{each} Chilean cohort. Across the four, adding the course in its reported position gives 846 fewer first choices, while putting it last gives 2,163 more. The option sets and the $+676$ historical-offer coverage difference are identical across these two scenarios. In the in-order scenario the blind rule gives both more placements and more first choices. It therefore wins for every nonnegative first-choice bonus in the index from Section~\ref{sec:ranks}. A returning pupil cannot reach an option below their protected course in the replay. Inserting the extra option immediately before that entitlement instead of after it gives $-1,764$ placements and $+2,161$ first choices across the cohorts.

The ordinary NYC first-choice differences for in-order restoration and last-place insertion are $-0.13$ and $+2.17$ points. The former interval spans zero ($[-0.47,0.22]$), as does the corresponding enriched interval ($[-0.16,0.54]$). If everyone tries all omitted options after the shortlist, the enriched placement difference estimates $+0.28$ points with a paired simulation-block interval of $[-0.07,0.65]$. The ordinary estimate is $-0.38$ with interval $[-0.66,-0.11]$. Access to the full set therefore substantially changes the housing result. If all applicants know their options and can rank them freely, truthful reporting is weakly dominant in the applicant-proposing DA model \citep{dubinsfreedman}. Restoring the full original order then makes the two rules identical. Trying the recommendations first is a behavioral assumption about how families use advice.

We also allowed each applicant, independently, to pursue their entire original list with probability one-half and otherwise use the five shown. The same draw applies to the person under both policies. The Chilean difference averaged 887 fewer placements (Monte Carlo SE 1.69) and 1,057 more first choices (SE 2.46). The NYC placement differences were $-1.03$ and $-1.63$ points for ordinary and enriched inventory. These probabilities describe a sensitivity analysis. Actual decisions to keep looking could depend on an applicant's preferences, information and expected chances.
